\section{Random-cluster conventions and imported estimates}
\label{sec:inputs}

We first identify the random couplings at the physical critical point.
We then state the two inputs from the companion paper~\cite{Q} that will
be used throughout the proof: a local bound on FK arms and a scale
transformation whose bulk perturbations vanish.  The transformation acts
on spin partition functions.  Its extension to the boundary of a polygon
is proved below; that extension is not part of the imported bulk estimate.

\subsection{The two probability spaces and planar duality}

For a fixed bond environment $J$, a finite graph $G=(V,E)$, and a
partition $\pi$ of its boundary vertices, write
\begin{equation}
 \phi_{G,J}^{\pi}(\omega)
 =\frac{1}{Z_{G,J}^{\pi}}\,
 2^{k_\pi(\omega)}\prod_{e\in E}v_e^{\omega_e},
 \qquad v_e=e^{2K_e}-1,\quad K_e=\beta J_e.
 \label{eq:fk-odds}
\end{equation}
Here $\omega\in\{0,1\}^E$, and $k_\pi$ counts open components after the
vertices in each part of $\pi$ have been identified.  A connection event
in this paper always uses actual open edges, unless a specified boundary
arc is explicitly allowed as part of the path.  In particular, identifying
two sets of vertices does not make an ordinary radial annulus crossing
occur automatically.

At fixed $J$, the law~\eqref{eq:fk-odds} has the spatial Markov property,
satisfies FKG, and is increasing in the odds and in the boundary
partition~\cite{FK,FKG}.  These statements remain valid for
inhomogeneous positive odds.  Planar FK duality gives
\begin{equation}
 v_{e^*}^{\,*}=\frac{2}{v_e}.
 \label{eq:dual-odds}
\end{equation}
When completing a subgraph to a larger graph for an upper bound, we
retain the actual odds on every shared edge and compare the induced
boundary partitions.

The Edwards--Sokal coupling~\cite{ES} relates this law to the Ising
weight $\exp(\sum_e K_e\sigma_x\sigma_y)$.  Given the spins, edges
joining equal spins are independently open with probabilities
$1-e^{-2K_e}$, and edges joining opposite spins are closed.  Given the
edges, each unwired component has an independent fair spin; a component
with a prescribed pin has that pin's spin.  Two separately wired arcs
can instead each carry one unconstrained terminal spin.  If
$Z^{++}$ and $Z^{+-}$ are the spin partition functions obtained by fixing
their two spins as indicated, while all other boundary vertices are
free, then their actual FK connection probability is
\begin{equation}
 \phi(\text{first arc}\leftrightarrow\text{second arc})
 =\frac{Z^{++}-Z^{+-}}{Z^{++}+Z^{+-}}.
 \label{eq:es-partition-ratio}
\end{equation}
Indeed the connection probability equals the correlation of the two
terminal spins, and global spin flip identifies the other two choices
of terminal signs.

We denote probability and expectation over the original iid signs by
$\P_{\mathrm{env}}$ and $\E_{\mathrm{env}}$.  The auxiliary priorities
and thresholds used to construct a scale transformation enlarge this
environment; $\E$ without a subscript in a moment bound includes that
construction randomness.  In contrast, independent \emph{reference
variables}, including the fair comparison coins used by the
transformation, belong to the Gibbs field: they are arguments of the
logarithmic weights and are sampled with that field.  Let
$K_0=\frac12\log(1+\sqrt2)$.  Pure expectation $\E_0$ is expectation
under the full-plane critical Ising law at coupling $K_0$, together with
the independent reference variables; it does not average the environment.
Supremum norms below range over spins and reference variables alike.
The original FK law and all its probabilities are independent of the
additional construction choices.

\subsection{The physical critical point and local barriers}

The coordinate used in~\cite{Q} is
\begin{equation}
 K(t)=\frac12\log(1+\sqrt2 e^t),\qquad
 K_e=K(h+\theta\xi_e).
 \label{eq:critical-coordinate}
\end{equation}
Thus $h$ changes both log odds by the same amount, while $\theta$ is
the size of their symmetric disorder.

\begin{proposition}[Critical-temperature input]\label{prop:critical-input}
For sufficiently small $\eps>0$, the plus-magnetization threshold for
$J_e=1+\eps\xi_e$ is the unique positive solution of
\begin{equation}
 \bigl(e^{2\beta_c(\eps)(1+\eps)}-1\bigr)
 \bigl(e^{2\beta_c(\eps)(1-\eps)}-1\bigr)=2.
 \label{eq:critical-self-duality}
\end{equation}
The plus magnetization vanishes at this temperature.  If
$v_\pm=e^{2\beta_c(\eps)(1\pm\eps)}-1$ and
$\theta=\frac12\log(v_+/v_-)$, then
\[
 v_e=\sqrt2\,e^{\theta\xi_e},\qquad \theta\longrightarrow0
 \quad\text{as }\eps\downarrow0.
\]
The scale trajectory with vanishing perturbations constructed in
\cite{Q} occurs at $h=0$ in~\eqref{eq:critical-coordinate}.
\end{proposition}

\begin{proof}[Source and normalization]
These are Proposition~8.2, \emph{Identification of the tuned parameter},
and Theorem~8.5, \emph{Physical critical point}, in Section~8 of~\cite{Q}.
The definition of magnetization there first normalizes each quenched
Gibbs law and then averages its infinite-volume plus magnetization.
It is therefore the threshold used here.  The bounds
$K_0/(1+\eps)\le\beta_c(\eps)\le K_0/(1-\eps)$ imply
$\theta\to0$.  The companion's identification of $h=0$ uses torus
partition functions and a separate sharpness argument to identify the
physical threshold.  No inference from self-duality alone is needed.
\end{proof}

Consequently the dual odds at this temperature have the same law as
the primal odds: the crossing-edge bijection in~\eqref{eq:dual-odds}
replaces $\xi_e$ by $-\xi_e$.  Statements in environmental probability
can therefore be applied to either color.

\begin{lemma}[Local FK barrier input]\label{lem:barrier-input}
There are fixed $a_0>1$ and $c_0>0$ such that the following holds at
$\beta_c(\eps)$.  Let
$B_r(x)=\{y\in\mathbb Z^2:\norm{y-x}_\infty\le r\}$, and let
$\mathcal A(x,r)=B_{a_0r}(x)\setminus B_r(x)$ be its lattice annulus.
With $\P_{\mathrm{env}}$-probability tending to one as $r\to\infty$,
\begin{equation}
 \phi_{\mathcal A(x,r),J}^{\mathrm w}
 (\text{ordinary open radial crossing})\le1-c_0.
 \label{eq:barrier-input}
\end{equation}
Here both rims are wired together.  The convergence is uniform in $x$,
and the event that~\eqref{eq:barrier-input} holds depends only on the
bonds of the annulus and its incident lattice layer.  The same bound
holds for every less wired boundary condition.  The assertion also
holds for dual crossings with dual odds.
\end{lemma}

\begin{proof}[Source]
This is Lemma~8.3, \emph{Local FK barriers}, in Section~8 of~\cite{Q},
followed by~\eqref{eq:dual-odds}.  That lemma gives no rate for the
exceptional environmental probability.  Our uses fix all physical
annular radii before taking the mesh limit.
\end{proof}

\subsection{Vanishing bulk perturbations}

The comparison of polygon partition functions will use the same local
transformation at many scales.  We record both its probabilistic output
and the precise exactness that survives a change of boundary geometry.
Fix an integer initial spacing $s_0$.  Allocate each lattice bond to its
nearest centers of the grid $s_0\mathbb Z^2$, splitting ties equally,
and denote the allocation weights by $\alpha_{i,e}$.  They are
nonnegative and sum to one over $i$ for each $e=\{x,y\}$.  Set
\begin{equation}
 f_i^{(0)}=\sum_e\alpha_{i,e}(K_e-K_0)b_e,
 \qquad b_e=\sigma_x\sigma_y-\E_0(\sigma_x\sigma_y).
 \label{eq:initial-logs}
\end{equation}
Their sum is the change in spin log weight, up to a scalar independent
of the spins.  The allocation is invariant under grid translations
and square symmetries.

\begin{proposition}[Bulk scale-transformation input]
\label{prop:bulk-input}
There is a fixed support constant $C_{\rm s}$ with the following
properties.  Choose a sufficiently large odd step factor $L$, a
sufficiently small flag threshold parameter $b_0>0$ thereafter, and a
sufficiently large fixed $s_0$.  For all sufficiently small $\theta>0$,
the plane fields~\eqref{eq:initial-logs} at $h=0$ admit successive
transformations to potentials $f_i^{(k)}$ at spacings $s_k=s_0L^k$ such
that:
\begin{enumerate}
\item Each $f_i^{(k)}$ is a bounded real spin-flip-even function, is
centered by $\E_0 f_i^{(k)}=0$, and has field support and environmental
dependence in the square of radius $C_{\rm s}s_k$ about $s_ki$.
The primitive environmental variables are independent.  Collections
with disjoint dependence neighborhoods are independent as random
functions.  The plane field is stationary on its grid and square
symmetric in distribution.
\item With the fixed moment order $P=40$,
\begin{equation}
 X_k=\bigl(\E\norm{f_i^{(k)}}_\infty^P\bigr)^{1/P}
 \longrightarrow0.
 \label{eq:bulk-decay}
\end{equation}
The numbers $X_k$ are uniformly arbitrarily small when $\theta$ is
sufficiently small.  The disorder bound and construction parameters
are chosen independently of any domain or mesh approximation.
\item Every local replacement preserves exterior-integrated weights
exactly, up to field-independent centering factors.  This holds for
arbitrary assignments of the old exterior field whenever the complete
local construction is performed.  For numerical $K_{\rm s}$ and a
constant $C_L$, one scale step satisfies
\begin{equation}
 \norm{f'_j}_\infty\le
 C_L\sum_{i:\,\norm{si-Lsj}_\infty\le K_{\rm s}Ls}
          \norm{f_i}_\infty.
 \label{eq:map-raw-input}
\end{equation}
The formulas are local and agree in embedded plane and torus charts.
\end{enumerate}
Any additional fixed lower bound on $L$, upper bound on $b_0$, or lower
bound on $s_0$ can be imposed in this order, followed by decreasing the
allowed disorder interval.
\end{proposition}

\begin{proof}[Source and parameter choices]
The local construction and~\eqref{eq:map-raw-input} are Proposition~3.3,
\emph{The local scale map}, of~\cite{Q}.  Proposition~4.1,
\emph{One-step estimates}, first makes the contraction coefficients
small by increasing $L$, then controls the Taylor remainder by decreasing
the flag threshold, and finally restricts the input moments.
Theorem~6.2, \emph{A small critical trajectory}, gives
$X_k\le A\sqrt{v_k}$ and
$v_k\le(v_0^{-1}+ck)^{-1}$ for fixed positive constants $A,c$ and
$v_0=O(\theta^2)$.  Here $v_k$ is a scalar variance coordinate of that
construction, distinct from the edge odds $v_e$.
Proposition~\ref{prop:critical-input} identifies the resulting trajectory
with $h=0$.  These proofs give the stated freedom in the parameter
choices; they do not require summability of $X_k$ over $k$.
\end{proof}

The exactness assertion concerns arbitrary assignments of the \emph{old}
exterior field, while the new Gibbs reference variables are sampled by
the complete local construction.  It does not permit those new variables
to be pinned separately.  Corollary~3.4, \emph{Local implementation},
of~\cite{Q} omits operations whose determining neighborhoods leave a
free chart and agrees with the plane construction only in an eroded
chart.  Section~\ref{sec:polygon} gives the construction, retains its
complete operations in a pad around the polygon, and proves the
additional boundary estimates.

There is also a distinction between local and global environmental
independence when using several sheets over the plane.  Different copies
of one physical edge carry the same bond environment.  We use a plane
moment estimate only when the whole determining neighborhood projects
injectively to an ordinary lattice chart.  Finite sums of final norms
are bounded by their individual moments; no independence between
separate sheets is required.
